% titel: Blatt 8, A2: Wasser höherer Siedepunkt als H2S
% skript: Mortimer Kap. 12 (Wasserstoffbrücken)

% Teil I: Werkzeuge, kein Lösungsweg.

Der Vergleich zweiter gleichartiger Verbindungen (H$_2$O gegen H$_2$S)
zeigt, wie eine besondere Kraftwirkung den ganz allgemeinen Trend
übersteigen kann.

\begin{werkzeugkasten}[Warum H$_2$O so hoch siedet]
Wasserstoffbrücken sind eine besonders starke spezifische intermolekulare
Anziehung zwischen dem H-Atom einer polaren Bindung und einem
freien Elektronenpaar eines elektronegativen Atoms (N, O, F).
Jedes H$_2$O-Molekül kann durch seine zwei H-Atome und zwei freie
Elektronenpaare ein räumliches Netzwerk von H-Brücken ausbilden.
\\[2mm]
\textit{Skript:} Mortimer Kap. 12 (Wasser), Kap. 11.
\end{werkzeugkasten}

\begin{achtung}[Darauf achten]
\begin{itemize}
  \item Einen abweichenden Gipfel in einer sonst steilen Reihe erklärt man mit
        einer \emph{extra starken} Wechselwirkung, nicht mit der Molmasse.
  \item Schwefel ist weniger elektronegativ als Sauerstoff $\Rightarrow$
        H$_2$S bildet fast keine H-Brücken.
  \item ``Was spräche für das Gegenteil?'' Der Aufgabe liegt eine Reihe
        zugrunde, die ohne H-Brücken nur nach der Molmasse liefe.
\end{itemize}
\end{achtung}
